%\input ../util/bookmacros
%\input ../util/docparams

\appendix{Appendix G}{Answers to the Exercises}
A word before you look. The appendix ``How to Study from This Book''
argues that an answer you fought for sticks far better than one you were
shown. So use these pages the way you would use an examiner: {\it commit\/}
to your own answer first -- in writing -- and only then compare. If your
answer differs, find the step where the two part ways before reading
on; that step is the most valuable thing these pages can give you.

The answers are brief. They show the move that was played and the result,
not every line of algebra.

\def\ans#1{\medskip\noindent{\bf Exercise #1.}\ }

\section{Difference Calculus}

\ans{1.1} Factor out $2^n$:
$\Delta 2^n = 2^{n+1} - 2^n = 2^n (2 - 1) = 2^n$.

\ans{1.2} The difference sequence is $\{1,\,3,\,5,\,7,\,9\}$ and the second
difference is $\{2,\,2,\,2,\,2\}$ -- a constant sequence. The original
sequence is $n^2$ (starting at $n = 0$), so this is $\Delta n^2 = 2n + 1$
followed by $\Delta (2n + 1) = 2$. Each differencing lowers the degree of a
polynomial by one, so two of them reduce a quadratic to a constant.

\ans{1.3} Group the product as $x_n (y_n z_n)$ and use the Discrete Product
Rule twice:
$$
\eqalign{
\Delta (x_n y_n z_n) & = \Delta x_n \, (y_{n+1} z_{n+1}) + x_n \, \Delta (y_n z_n) \cr
& = \Delta x_n \, y_{n+1} z_{n+1} + x_n \, \Delta y_n \, z_{n+1} + x_n y_n \, \Delta z_n \cr
}
$$
Each term differences exactly one factor; the factors to its right are
evaluated at $n+1$ and those to its left at $n$. (Other groupings give
other, equally correct, forms.)

\section{Summation Calculus}

\ans{2.1} By the falling power rule $\Delta n^{\underline{3}} = 3\, n^{\underline{2}}$,
so an anti-difference of $n^{\underline{2}}$ is $n^{\underline{3}}/3$. By the FTDC
$$
\sum_{n=1}^N n(n-1) = {(N+1)^{\underline{3}} \over 3} - {1^{\underline{3}} \over 3}
= {(N+1)\,N\,(N-1) \over 3}
$$
since $1^{\underline{3}} = 1 * 0 * (-1) = 0$. Check: for $N = 3$ the sum is
$0 + 2 + 6 = 8$ and the formula gives $4 * 3 * 2 / 3 = 8$.

\ans{2.2} Use Summation by Parts with $x_n = n^3$ and $y_n = n$, exactly as in
the example for $\sum n^2$. Since $\Delta n^3 \, (n+1) = (3n^2 + 3n + 1)(n+1) =
3n^3 + 6n^2 + 4n + 1$, we get
$$
4 \sum_{n=1}^N n^3 = (N+1)^4 - 1 - 6 \sum_{n=1}^N n^2 - 4 \sum_{n=1}^N n - N
$$
Substituting the formulas already found for $\sum n^2$ and $\sum n$ and
simplifying, the right hand side becomes $N^2 (N+1)^2$. Therefore
$$
\sum_{n=1}^N n^3 = {N^2 (N+1)^2 \over 4} = \biggl( {N(N+1) \over 2} \biggr)^2
$$
-- the square of the sum of the first $N$ integers. Check: for $N = 2$,
$1 + 8 = 9 = 3^2$.

\ans{2.3} With $x_n = n^2$ and $y_n = 2^n$ we have $\Delta x_n = 2n + 1$ and
$\Delta y_n = 2^n$, so Summation by Parts gives
$$
\sum_{n=1}^N n^2 \, 2^n = (N+1)^2 \, 2^{N+1} - 2 - \sum_{n=1}^N (2n+1)\, 2^{n+1}
$$
The remaining sum is $4 \sum n\, 2^n + \sum 2^{n+1} =
4\bigl((N-1)\,2^{N+1} + 2\bigr) + \bigl(2^{N+2} - 4\bigr) = (4N - 2)\, 2^{N+1} + 4$.
Therefore
$$
\sum_{n=1}^N n^2 \, 2^n = (N^2 - 2N + 3)\, 2^{N+1} - 6
$$
Check: for $N = 2$ the sum is $2 + 16 = 18$ and the formula gives
$3 * 8 - 6 = 18$.

\section{Differential Calculus}

\ans{3.1} For a sequence $h_n \ne 0$ tending to $0$, expand the cube:
$$
{(x + h_n)^3 - x^3 \over h_n} = {3x^2 h_n + 3x h_n^2 + h_n^3 \over h_n}
= 3x^2 + 3x\, h_n + h_n^2
$$
The last two terms tend to $0$, so $f'(x) = 3x^2$. (Compare
$\Delta n^3 = 3n^2 + 3n + 1$: the same expression with step size $1$.)

\ans{3.2} Let $x$ be the length of each side perpendicular to the river; the
side along the river is then $100 - 2x$, and $x$ lies in $[0, 50]$. The
area is $A(x) = x(100 - 2x) = 100x - 2x^2$, so $A'(x) = 100 - 4x$, which
is $0$ only at $x = 25$. The candidates are $\{0, 25, 50\}$ with
$A(0) = A(50) = 0$ and $A(25) = 1250$. The best pen is $25$ feet by $50$
feet (the long side along the river), enclosing $1250$ square feet.

\ans{3.3} (1) Chain rule, with outer function $z^{10}$ and inner function
$1 + x^3$: $f'(x) = 10 (1 + x^3)^9 * 3x^2 = 30\, x^2 (1 + x^3)^9$.

\noindent (2) Product rule, with the chain rule for the second factor:
$f'(x) = 2x (1+x)^5 + 5x^2 (1+x)^4 = x\,(1+x)^4 (2 + 7x)$.

\noindent (3) Quotient rule:
$f'(x) = {1 * (1 + x^2) - x * 2x \over (1 + x^2)^2} = {1 - x^2 \over (1 + x^2)^2}$.

\section{Integral Calculus}

\ans{4.1} Let $H = F - G$. By linearity $H' = F' - G' = f - f = 0$, so the
graph of $H$ has a flat tangent at every point. A function whose slope is
$0$ everywhere never rises or falls: it is a constant function. (The Mean
Value Theorem of the chapter ``The Derivative As Approximation'' makes
this airtight: $H(b) - H(a) = H'(\mu)(b - a) = 0$ for any $a$ and $b$.)
So $F(x) - G(x) = c$; that is, $F(x) = G(x) + c$.

\ans{4.2} (1) Substitution with $u(x) = 1 + x$, $u'(x) = 1$:
$\int_0^1 (1+x)^7 \, dx = {(1+x)^8 \over 8}\,\big|_{x=1} - {(1+x)^8 \over 8}\,\big|_{x=0}
= {2^8 - 1 \over 8} = {255 \over 8}$.

\noindent (2) Integration by parts with $f(x) = x$ and $g(x) = \sin(x)$, so
that $g'(x) = \cos(x)$:
$\int_0^{\pi/2} x \cos(x) \, dx = \bigl( (\pi/2) \sin(\pi/2) - 0 \bigr) -
\int_0^{\pi/2} \sin(x) \, dx = \pi/2 - 1$.

\noindent (3) Substitution with $u(x) = x^2$, $u'(x) = 2x$; the integrand is
${1 \over 2} e^{u(x)} u'(x)$:
$\int_0^1 x\, e^{x^2} \, dx = {1 \over 2} e^{x^2}\,\big|_{x=1} - {1 \over 2} e^{x^2}\,\big|_{x=0}
= {e - 1 \over 2}$.

\ans{4.3} An anti-derivative of $\sin$ is $-\cos$, so
$\int_0^{\pi} \sin(x) \, dx = -\cos(\pi) + \cos(0) = 1 + 1 = 2$, while
$\int_0^{2\pi} \sin(x) \, dx = -\cos(2\pi) + \cos(0) = -1 + 1 = 0$. On
$[\pi, 2\pi]$ the sine is negative, so the terms $f(t_n) \Delta t_n$ of the
approximating sums are negative there: the integral counts area below the
axis as negative, and the two humps cancel. The area actually enclosed is
$2 + 2 = 4$.

\section{The Derivative As Approximation}

\ans{5.1} With $f(x) = x^5$ we have $f'(x) = 5x^4$, so $f(2) = 32$ and
$f'(2) = 80$. Then $(2.01)^5 \approx 32 + 80 * 0.01 = 32.8$. The estimate
is low by about $0.008$ -- an error of roughly $0.025$\%.

\ans{5.2} The estimate is $|V'(60) * 2| = 14{,}400 * 2 = 28{,}800$ cubic
inches, about $0.62$ cubic yards. Yes, this is exactly twice the $1$ inch
estimate: the tangent line approximation is {\it linear\/} in $h$. The
exact extra amount is $V(62) - V(60) = 120\,(62^2 - 60^2) = 29{,}280$ cubic
inches, which is {\it not\/} twice the exact $1$ inch amount
($2 * 14{,}520 = 29{,}040$). The estimate is short by $480 = 120 * 2^2$
cubic inches -- the error grows like $h^2$, as Taylor's Theorem predicts.

\ans{5.3} With $f = \cos$ and $c = 0$: $f(0) = 1$, $f'(0) = -\sin(0) = 0$,
$f''(0) = -\cos(0) = -1$, and $f''' = \sin$. Taylor's Theorem with
$n = 3$ gives
$$
\cos(x) = 1 - {x^2 \over 2} + \sin(\mu(x)) {x^3 \over 3!}
$$
Since the sine is never bigger than $1$ in size, the error is at most
$|x|^3/6$. For $|x| \le 0.1$ that is at most $(0.1)^3/6 \approx 1.7 \times
10^{-4}$.

\section{Difference Equations}

\ans{6.1} Here $a = 2$, $b = 1$, and $L = 0$, so
$x_N = a^N L + b\,{1 - a^N \over 1 - a} = 0 + {1 - 2^N \over 1 - 2} = 2^N - 1$.

\ans{6.2} The balance grows by the interest and by the deposit:
$x_{n+1} = (1 + r)\, x_n + p$, with $x_0 = L$. This is the equation of the
chapter with $a = 1 + r$ and $b = p$, so
$$
x_N = (1+r)^N L + p \biggl( {(1+r)^N - 1 \over r} \biggr)
$$
It is the loan formula with the sign of the second term reversed: in the
loan the payments work against the compounding; in the savings account
the deposits work with it.

\ans{6.3} Write $\lambda_1 = (1 + \sqrt 5)/2$ and $\lambda_2 = (1 - \sqrt 5)/2$;
note that $\lambda_1 + \lambda_2 = 1$ and $\lambda_1 - \lambda_2 = \sqrt 5$.
The initial values require $A + B = 1$ and $A \lambda_1 + B \lambda_2 = 1$.
Solving, $A = \lambda_1 / \sqrt 5$ and $B = -\lambda_2 / \sqrt 5$. Therefore
$$
F_n = {\lambda_1^{n+1} - \lambda_2^{n+1} \over \sqrt 5}
$$
Check: since each $\lambda$ satisfies $\lambda^2 = \lambda + 1$, we have
$\lambda^3 = 2 \lambda + 1$ and $\lambda^4 = 3 \lambda + 2$. So
$F_2 = 2 (\lambda_1 - \lambda_2)/\sqrt 5 = 2$ and
$F_3 = 3 (\lambda_1 - \lambda_2)/\sqrt 5 = 3$, in agreement with
$1, 1, 2, 3, \ldots$

\section{Differential Equations}

\ans{7.1} Group the unknowns: $T'(t) + k\, T(t) = k A$. Multiplying by $y(t)$
and matching against $(y\, T)' = y\, T' + y'\, T$ requires $y' = k\, y$; so
the integrating factor is $e^{k t}$, and
$\bigl( e^{k t}\, T(t) \bigr)' = k A\, e^{k t}$. Integrating from $0$ to $t$
gives $e^{k t}\, T(t) - T_0 = A\, (e^{k t} - 1)$; that is,
$$
T(t) = A + (T_0 - A)\, e^{-k t}
$$
As $t$ grows large $e^{-k t}$ tends to $0$, so the coffee's temperature
tends to the room temperature $A$ -- quickly at first, then ever more
slowly.

\ans{7.2} The amount is $a(t) = Q\, e^{-r t}$. The half-life gives
$e^{-100\, r} = {1 \over 2}$, so $r = \ln(2)/100$. We need
$e^{-r t} = {1 \over 10}$, so
$t = \ln(10)/r = 100 \ln(10)/\ln(2) \approx 332$ years -- a little more
than three half-lives, as it should be, since three half-lives leave
$1/8 = 12.5$\%.

\ans{7.3} $(1.005)^{60} \approx 1.3489$ while $e^{0.3} \approx 1.3499$. They
differ by about $0.001$ -- less than a tenth of a percent.

\ans{7.4} Differentiating twice: $\sin' = \cos$ and $\cos' = -\sin$, so
$\sin'' = -\sin$; likewise $(-\sin)' = -\cos$, so
$\cos'' = -\cos$. In both cases $x'' + x = 0$. For the combination
$x(t) = A \sin(t) + B \cos(t)$, linearity of differentiation gives
$x''(t) = -A \sin(t) - B \cos(t) = -x(t)$, so again $x'' + x = 0$.

\section{The Exercise Left to the Reader in Appendix D}
The appendix on Mathematical Induction left one claim unproved:
$(x^{-1})' = -x^{-2}$ for $x \ne 0$. From the definition, for a sequence
$h_n \ne 0$ tending to $0$ (and far enough out that $x + h_n \ne 0$),
$$
{1 \over h_n} \biggl( {1 \over x + h_n} - {1 \over x} \biggr)
= {1 \over h_n} \biggl( {x - (x + h_n) \over x\,(x + h_n)} \biggr)
= {-1 \over x\,(x + h_n)}
$$
which tends to $-1/x^2$ as $h_n$ tends to $0$.

\vfil
